% Document E: Riccati Flow of the Hessian under Heat-Flow RG
% Suggested file name: riccati_flow_hessian.tex
% =========================================================

\section{Riccati Flow of the Hessian under Heat-Flow Coarse-Graining}

Let $S_t(x)$ be a family of effective actions on $\mathbb{R}^d$ defined via the heat equation
on the density:
\[
\partial_t \rho_t = \Delta \rho_t, \qquad \rho_t(x) = e^{-S_t(x)}.
\]

\subsection{Viscous Hamilton--Jacobi equation}

Substituting $\rho_t = e^{-S_t}$ gives
\[
\partial_t \rho_t = -(\partial_t S_t) e^{-S_t}, \quad
\Delta \rho_t = e^{-S_t}(-\Delta S_t + |\nabla S_t|^2).
\]
Thus
\[
-(\partial_t S_t)e^{-S_t}
= e^{-S_t}(-\Delta S_t + |\nabla S_t|^2),
\]
and hence
\[
\partial_t S_t = \Delta S_t - |\nabla S_t|^2.
\]

\subsection{Equation for the Hessian}

Let $H_t(x) = \nabla^2 S_t(x)$ be the Hessian matrix with entries
\[
(H_t)_{ij}(x) = \partial_i\partial_j S_t(x).
\]

Differentiating $\partial_t S_t = \Delta S_t - |\nabla S_t|^2$ twice,
\[
\partial_t(\partial_i\partial_j S_t)
= \partial_i\partial_j \Delta S_t
 - \partial_i\partial_j(|\nabla S_t|^2).
\]

The first term is
\[
\partial_i\partial_j \Delta S_t = \Delta(\partial_i\partial_j S_t) = (\Delta H_t)_{ij}.
\]

For the second term, write $|\nabla S_t|^2 = \sum_k (\partial_k S_t)^2$ and compute
\[
\partial_i\partial_j |\nabla S_t|^2
= 2\sum_k \big( (\partial_i\partial_k S_t)(\partial_j\partial_k S_t)
 + (\partial_k S_t)(\partial_i\partial_j\partial_k S_t)\big).
\]

Hence
\[
\partial_t H_{ij}
= (\Delta H_t)_{ij}
 - 2\sum_k (\partial_i\partial_k S_t)(\partial_j\partial_k S_t)
 - 2\sum_k (\partial_k S_t)(\partial_i\partial_j\partial_k S_t).
\]

In matrix form this is
\[
\partial_t H_t
= \Delta H_t - 2H_t^2 + R_t,
\]
where the remainder $R_t$ has components
\[
(R_t)_{ij}
:= -2\sum_k (\partial_k S_t)(\partial_i\partial_j\partial_k S_t).
\]

Thus the Hessian of the effective action under heat-flow coarse-graining
obeys a matrix Riccati-type equation with a Laplacian term and a remainder.
This is the flat-space analogue of the Riccati flow invoked in the continuum conjecture.%
\footnote{Cf.\ Section~6.1 of the monograph. :contentReference[oaicite:2]{index=2}}
